\section{Introdução}
\label{sec:introducao}

Este documento serve como exemplo de projeto no editor, reunindo os
elementos mais comuns de um texto acadêmico: equações numeradas, tabelas,
figuras, múltiplos arquivos \texttt{.tex} e uma bibliografia em
\texttt{.bib}.

Um exemplo clássico é o modelo de regressão linear simples, que relaciona
uma variável resposta $y$ a uma variável explicativa $x$:
\begin{equation}
    y = \beta_0 + \beta_1 x + \varepsilon,
    \label{eq:regressao}
\end{equation}
em que $\beta_0$ é o intercepto, $\beta_1$ é o coeficiente angular e
$\varepsilon$ é o erro aleatório \cite{montgomery2012}. A
Equação~\eqref{eq:regressao} será usada como referência ao longo deste
exemplo.

Também é possível citar outros trabalhos, como o artigo clássico de
Einstein sobre eletrodinâmica \cite{einstein1905} ou o trabalho de
Cortes e Vapnik sobre máquinas de vetores de suporte \cite{cortes1995}.
